% =====================================================================
% Chapter 7 -- Experiment 3: Reconstruction-Augmented Forecasting
%              Without Imputation -- Results and Discussion
% ---------------------------------------------------------------------
% Required packages (preamble): booktabs, multirow, amsmath, graphicx
% Lines marked "% AUTHOR-CHECK" are claims that depend on implementation
% details the author should confirm before submission.
% =====================================================================

\section{Results}
\label{sec:exp3-results}

This section reports the empirical behaviour of RIoT v4 on the held-out
chronological test partition. It begins by setting out exactly what the
reported metrics measure, because several of the conclusions drawn later
depend on how missing observations enter the evaluation targets. It then
places the v4 results against the historical v3 record, examines the
contribution of spatial message passing relative to a non-spatial
baseline, and deconstructs the architecture through four single-factor
ablations and two optimisation variants.

\subsection{Evaluation protocol and metric definitions}
\label{sec:exp3-protocol}

All v4 configurations were trained and evaluated under the protocol
introduced in Section~\ref{sec:exp3-method}: a strict chronological
70/15/15 split across the multi-month record, a purge gap of
$T + h = 13$ steps (3.25\,h) between consecutive partitions, and
standardisation statistics estimated on the training partition alone.
The network comprises all 81 sensor locations, of which 24 are designated
as inductive stations. These stations are excluded from every
gradient-producing loss term during training and are encountered only at
evaluation time, so that kriging performance on them constitutes a
genuine zero-shot spatial transfer test. The inductive set was drawn
once, with a fixed seed (42), and is shared across all configurations,
so every run is scored on the same 24 stations.

The forecast metrics (RMSE, MAE and ZIR) are averaged over \emph{all
81 nodes}, not over the 57 transductive nodes alone. The inductive
stations are excluded from the training losses, and at evaluation time
their dynamic inputs are zeroed and their total-missingness channel is
set to ``unobserved''. They nonetheless still receive horizon-$h$
forecasts, and those forecasts are scored alongside the transductive
ones. Each headline forecast figure is therefore a blend: roughly
70\% of it reflects transductive forecasting from a node's own history,
and roughly 30\% reflects zero-shot forecasting at stations whose
history the model never sees. This blend affects every configuration
equally, so comparisons between runs remain valid. Absolute values,
however, should not be read as transductive forecast skill.
Section~\ref{sec:exp3-limitations} lists the stratified re-scoring this
calls for.

One implementation detail also qualifies the zero-shot claim. When the
inductive stations are blanked, only the total-missingness channel is
overwritten. The night-restricted missingness channel
(Section~\ref{sec:exp3-zero-recall}) keeps its true values. The model
can therefore see exactly when each inductive station was unobserved
at night during the input window, and because targets are zero-filled,
this reveals when that station's targets are likely to be zero. The
leak is narrow, since it carries no count values, daytime outage
pattern or observed-zero information. It does, however, bear directly
on the kriging metrics, on forecast error at the inductive stations,
and on ZIR. The kriging results below should be read as provisional
until they are re-scored with this channel blanked for held-out
stations.

Four properties of the metric implementation matter for interpretation.

\begin{enumerate}
    \item \textbf{Forecast RMSE and MAE} are computed on raw pedestrian
    counts at horizon $h$, with the point forecast taken as the ZINB
    expectation $\hat{y} = \hat{\mu}(1 - \hat{\pi})$.
    \item \textbf{Forecast targets at unobserved timestamps are zero.}
    Missing values in the count series are zero-filled before targets
    are constructed, and the ZINB likelihood is evaluated on those
    targets without a validity mask. An interval in which a sensor
    failed to report is therefore scored, and trained, as though zero
    pedestrians passed. The missingness mask channels remain the only
    signal by which the model can tell an observed zero from an
    unobserved interval.
    \item \textbf{The Zero-Identification Rate (ZIR)} is the proportion
    of target positions with $y = 0$ for which the zero-inflation head
    assigns $\hat{\pi} > 0.5$. The implementation logs this quantity as
    \texttt{zero\_acc}, but it is \emph{not} a measure of accuracy, and
    it should not be maximised. Given (2), its denominator combines
    genuinely observed zeros with \emph{all} zero-filled unobserved
    intervals, at any time of day. Only some of those intervals are
    plausible physical zeros. A ZIR of 100\% would mean the model
    treats every sensor outage as an absence of pedestrians, which is
    physically implausible for daytime gaps. We therefore read ZIR as
    the share of zero-valued targets that the model identifies as
    plausible structural zeros, and interpret it against a plausibility
    band (Section~\ref{sec:exp3-zero-recall}), not against a
    maximum. No corresponding precision is reported.
    \item \textbf{Imputation and kriging errors are reported in
    standardised ($z$-score) space.} Imputation MSE is evaluated on
    artificially masked entries across all dynamic channels. Kriging
    RMSE and MAE are computed across all dynamic channels (traffic and
    the two meteorological covariates) of the 24 inductive stations.
    Kriging MSLE is computed on the traffic channel only, after clamping
    standardised values at zero, which means it scores
    above-mean activity only. None of these three quantities is on the
    same scale as the forecast RMSE, and they should not be compared
    with it numerically.
\end{enumerate}

We call the design in (2) and (3) \emph{forecasting without imputation}
in the sense intended throughout this thesis: no statistical or learned
imputation is applied before modelling. It is not a model-free treatment
of missingness, however. Zero-filling combined with explicit mask
channels is itself a representational choice, and
Section~\ref{sec:exp3-zero-recall} shows that it has direct consequences
for what the zero-inflation head learns.

Each configuration is reported from a single training run.
% AUTHOR-CHECK: confirm single seed per configuration.
Differences of the order of 0.5--1.0 RMSE between configurations should
therefore be treated as indicative, not as established, until they are
replicated across seeds (Section~\ref{sec:exp3-limitations}).

% ---------------------------------------------------------------------
\subsection{Integrity versus performance: RIoT v3 and RIoT v4}
\label{sec:exp3-v3-v4}

Table~\ref{tab:exp3-v3-v4} sets out the methodological differences
between the two pipeline generations alongside their headline error
figures.

\begin{table}[htbp]
    \centering
    \caption[Methodological and empirical comparison of RIoT v3 and RIoT v4]%
    {Methodological and empirical comparison of the RIoT v3 and RIoT v4
    experimental pipelines. v3 figures are validation-set values under
    overlapping sliding-window splits; no out-of-sample test partition
    existed in v3. Error figures are not directly commensurable across
    generations (see text).}
    \label{tab:exp3-v3-v4}
    \small
    \begin{tabular}{@{}p{0.26\textwidth}p{0.32\textwidth}p{0.34\textwidth}@{}}
        \toprule
        \textbf{Dimension} & \textbf{RIoT v3} & \textbf{RIoT v4} \\
        \midrule
        \multicolumn{3}{@{}l}{\emph{Protocol}} \\
        Data split & Ad-hoc 24\,h slice or random sliding windows
                   & Chronological 70/15/15 \\
        Window overlap across folds & Up to 11 of 12 steps shared
                   & None; 13-step (3.25\,h) purge gap \\
        Scaler fitting & Global (full record) & Training partition only \\
        Test partition & None & Held-out chronological test set \\
        \midrule
        \multicolumn{3}{@{}l}{\emph{Spatial scope}} \\
        Real sensor nodes & 61 & 81 (incl.\ 27 reclaimed peripheral stations) \\
        Zero-history nodes & 5 simulated; degenerate ($r = 1.00$)
                   & 24 real inductive stations \\
        \midrule
        \multicolumn{3}{@{}l}{\emph{Model and objective}} \\
        Temporal positional encoding & No & Sinusoidal \\
        Input channels & Value only & Value, sensor status, artificial mask, MNAR mask \\
        Reconstruction loss & Raw-count MSE (mean 1{,}537; peaks $>$94{,}000)
                   & Normalised imputation loss \\
        Kriging supervision & No explicit inductive term
                   & $w_{\text{kriging}}\mathcal{L}_{\text{kriging}}$ \\
        \midrule
        \multicolumn{3}{@{}l}{\emph{Headline forecast error (raw counts)}} \\
        Validation RMSE & 21.51--26.56 (23 runs; best 21.51)
                   & 17.94 (full); 16.34 (bs\,=\,32) \\
        Validation MAE & 5.62--6.38 (representative runs)
                   & 5.58 (full); 5.88 (bs\,=\,32) \\
        Test RMSE / MAE & -- & 18.16 / 5.17 (full); 16.84 / 5.42 (bs\,=\,32) \\
        \bottomrule
    \end{tabular}
\end{table}
% AUTHOR-CHECK: the v3 "Input channels" and "Kriging supervision" rows are
% inferred from the chapter's description of v4 additions; verify against
% the v3 configuration.

The best v4 validation RMSE (17.94 for the reference configuration) is
16.6\% below the best v3 validation RMSE (21.51), and the v4 test RMSE
(18.16) is lower than every v3 validation figure. This is not the
direction one would expect after removing leakage. Overlapping windows
and global scaling both bias error estimates \emph{downwards}, so a
faithful re-evaluation of the v3 model under v4 hygiene would be expected
to produce higher errors than those originally reported, not lower ones.

We do not read the v3--v4 gap as evidence of an architectural gain of
roughly 16\%. At least three confounders separate the two sets of
numbers:

\begin{enumerate}
    \item \textbf{Different populations of nodes.} The 27 reclaimed
    peripheral stations are, almost by construction, lower-footfall
    locations. Because RMSE in raw counts scales with the magnitude of
    the series, adding low-volume nodes lowers network-averaged RMSE
    mechanically, even if per-node skill is unchanged.
    \item \textbf{Different temporal coverage.} The v4 test partition is
    the final 15\% of the record. Its seasonal and event composition
    differs from both the v3 24-hour slice and the randomly sampled v3
    windows.
    \item \textbf{Different target treatment.} Section~\ref{sec:exp3-zero-recall}
    shows that the v4 zero-inflation head can drive $\hat{y}$ towards
    zero at zero-filled outage intervals. This reduces error at
    positions that, in a physical sense, carry no ground truth at all.
\end{enumerate}

The defensible claim is therefore narrower. Under a protocol that is
strictly more conservative than v3's, RIoT v4 reaches forecast errors at
least as low as v3 reported under a leaky protocol. The v4 validation
estimates also transfer to the test partition with little degradation:
the reference model's RMSE rises from 17.94 to 18.16 (+1.2\%) between
validation and test, and its MAE falls from 5.58 to 5.17. A small,
bidirectional validation--test gap of this kind suggests that neither
model selection nor early stopping has overfitted the validation fold.
It is the main affirmative result of the protocol overhaul. It shows
that v4 figures can be trusted as forecasts of out-of-sample behaviour,
which v3 figures could not. A like-for-like attribution of the
architectural contribution would require re-scoring v4 on the original
61-node subset, or reporting a scale-normalised error such as RMSE
divided by per-node mean footfall. We leave this as an explicit item of
further work, and do not infer it from the present comparison.

% ---------------------------------------------------------------------
\subsection{Test-set performance of the v4 configurations}
\label{sec:exp3-main}

Table~\ref{tab:exp3-main} reports all eight v4 configurations on the
held-out test partition. Table~\ref{tab:exp3-delta} expresses each as a
relative change from the reference full model (\texttt{wph1nhvj}).

\begin{table}[htbp]
    \centering
    \caption[RIoT v4 test-set results]%
    {RIoT v4 results on the held-out chronological test partition
    (81 nodes; 24 inductive stations). Forecast RMSE and MAE are in raw
    counts. Imputation MSE and kriging metrics are in standardised space
    (Section~\ref{sec:exp3-protocol}). Best value per column in bold;
    arrows give the direction of improvement.}
    \label{tab:exp3-main}
    \resizebox{\textwidth}{!}{%
    \begin{tabular}{@{}llrrrrrrr@{}}
        \toprule
        & & \multicolumn{3}{c}{\textbf{Forecast}} & \textbf{Recon.}
          & \multicolumn{3}{c}{\textbf{Kriging (inductive)}} \\
        \cmidrule(lr){3-5} \cmidrule(lr){6-6} \cmidrule(l){7-9}
        \textbf{Configuration} & \textbf{Run}
          & RMSE$\downarrow$ & MAE$\downarrow$ & ZIR
          & Imp.\ MSE$\downarrow$
          & RMSE$\downarrow$ & MAE$\downarrow$ & MSLE$\downarrow$ \\
        \midrule
        Full model (reference)       & \texttt{wph1nhvj} & 18.164 & \textbf{5.174} & 0.554 & \textbf{0.0241} & 0.2351 & 0.0520 & \textbf{0.0356} \\
        Non-spatial LSTM             & \texttt{l5ck2j76} & 22.237 & 11.501 & --    & 0.1246 & 0.3873 & 0.2829 & 0.0385 \\
        \midrule
        \multicolumn{9}{@{}l}{\emph{Ablations}} \\
        \quad No imputation loss     & \texttt{vkk1ot7g} & 19.538 & 5.678 & 0.378 & 0.1521 & 0.2490 & 0.0568 & 0.0381 \\
        \quad No kriging loss        & \texttt{nfdgnlp9} & 18.754 & 5.298 & 0.602 & 0.0265 & 0.4397 & 0.3321 & 0.0794 \\
        \quad No missingness masks   & \texttt{ak47uk50} & 19.912 & 6.189 & 0.295 & 0.0305 & 0.2683 & 0.0734 & 0.0381 \\
        \quad No positional encoding & \texttt{pcj4st27} & 18.977 & 5.448 & 0.595 & 0.0265 & 0.2482 & 0.0554 & 0.0372 \\
        \midrule
        \multicolumn{9}{@{}l}{\emph{Optimisation variants}} \\
        \quad Batch size 32          & \texttt{iou8v4o8} & \textbf{16.839} & 5.416 & 0.386 & 0.0268 & \textbf{0.2133} & 0.0526 & 0.0374 \\
        \quad Dropout 0.3            & \texttt{9f0g2fbv} & 19.527 & 5.570 & 0.362 & 0.0276 & 0.2497 & \textbf{0.0511} & 0.0375 \\
        \bottomrule
    \end{tabular}%
    }
\end{table}

\begin{table}[htbp]
    \centering
    \caption[Relative change against the full model]%
    {Relative change of each configuration against the full model
    (\texttt{wph1nhvj}). Positive values indicate higher error, except
    for ZIR, which is given as an absolute change in percentage
    points (pp). Single-run estimates.}
    \label{tab:exp3-delta}
    \small
    \begin{tabular}{@{}lrrrrrrr@{}}
        \toprule
        & \multicolumn{3}{c}{\textbf{Forecast}} & \textbf{Recon.}
          & \multicolumn{3}{c}{\textbf{Kriging}} \\
        \cmidrule(lr){2-4} \cmidrule(lr){5-5} \cmidrule(l){6-8}
        \textbf{Configuration}
          & $\Delta$RMSE & $\Delta$MAE & $\Delta$ZIR (pp)
          & $\Delta$MSE
          & $\Delta$RMSE & $\Delta$MAE & $\Delta$MSLE \\
        \midrule
        Non-spatial LSTM       & +22.4\% & +122.3\% & --     & +417\% & +64.7\% & +444\% & +8.1\% \\
        No imputation loss     & +7.6\%  & +9.7\%   & $-17.6$ & +531\% & +5.9\%  & +9.2\% & +7.0\% \\
        No kriging loss        & +3.2\%  & +2.4\%   & $+4.8$  & +10.0\% & +87.0\% & +539\% & +123\% \\
        No missingness masks   & +9.6\%  & +19.6\%  & $-25.9$ & +26.6\% & +14.1\% & +41.2\% & +7.0\% \\
        No positional encoding & +4.5\%  & +5.3\%   & $+4.2$  & +10.0\% & +5.6\%  & +6.5\% & +4.5\% \\
        \midrule
        Batch size 32          & $-7.3$\% & +4.7\%  & $-16.8$ & +11.2\% & $-9.3$\% & +1.2\% & +5.1\% \\
        Dropout 0.3            & +7.5\%  & +7.7\%   & $-19.2$ & +14.5\% & +6.2\%  & $-1.7$\% & +5.3\% \\
        \bottomrule
    \end{tabular}
\end{table}

Three regularities are visible before any single mechanism is examined.
First, every ablation increases forecast RMSE and MAE relative to the
reference model. No component can be removed without some cost to
forecasting, although the cost ranges from marginal (kriging loss,
+3.2\% RMSE) to substantial (missingness masks, +19.6\% MAE). Second, the
ablations are largely \emph{selective}. Each one degrades most sharply
the output its component was designed to support: removing the kriging
loss damages kriging, removing the imputation loss damages
reconstruction, and removing the masks damages zero-classification.
This selectivity is what one would hope to see if the multi-task
objective is genuinely factorised and not simply over-parameterised.
Third, ZIR does not move with forecast accuracy. The two
configurations with the highest ZIR (no kriging loss, 0.602; no
positional encoding, 0.595) are both less accurate than the reference
model, while the most accurate configuration by RMSE (batch size 32) has
one of the lowest (0.386). This dissociation is taken up in
Section~\ref{sec:exp3-zero-recall}.

% ---------------------------------------------------------------------
\section{Discussion}
\label{sec:exp3-discussion}

\subsection{The inductive value of spatial representation}
\label{sec:exp3-spatial}

On every metric it reports, the non-spatial LSTM baseline
(\texttt{l5ck2j76}) is worse than the full model. Forecast RMSE is 22.4\%
higher, forecast MAE more than doubles (11.50 against 5.17), and on the
24 inductive stations kriging RMSE and MAE are 64.7\% and 444\% higher.
Taken at face value, this would be strong evidence that spatial message
passing contributes to both transductive forecasting and zero-shot
kriging. Two features of the comparison call for a more measured
reading.

The first is that the baseline differs from the full model in more than
its spatial component. It has no zero-inflation head, so its point
forecast is $\hat{\mu}$ and not $\hat{\mu}(1-\hat{\pi})$. It therefore
cannot zero out the forecast at the zero-filled outage intervals that
make up a substantial fraction of the targets
(Section~\ref{sec:exp3-zero-recall}). An error concentrated at positions
where $y=0$ inflates MAE disproportionately relative to RMSE, and this is
exactly the signature observed. The RMSE/MAE ratio is 1.93 for the LSTM,
against 3.51 for the full model, which indicates that the LSTM's error
is spread broadly across many moderate misses, not concentrated in a few
large ones. A substantial share of the 55\% MAE gap is therefore more
plausibly attributable to the distributional head than to the graph. The
comparison establishes that the \emph{full RIoT architecture} is
superior to a recurrent baseline. It does not isolate the contribution
of spatial message passing. The baseline also receives neither the
missingness-mask channels nor static node covariates. Its only input is
the zero-filled dynamic series, and it is trained with a plain MSE loss.
It therefore lacks both the mask-driven zero mechanism discussed in
Section~\ref{sec:exp3-zero-recall} and any means of telling an outage
from an observed zero.

The second concerns what the kriging metrics can distinguish. The LSTM
has no mechanism for propagating information from observed neighbours
to an unseen station. Its kriging output for an inductive station is
simply its reconstruction of that station's zeroed input, produced
without any location information. This makes it a strict ``spatially
blind'' reference point. The
full model's kriging MAE (0.052 in standardised units) is roughly
one-fifth of this reference (0.283), which is direct evidence that
information is being transferred from observed to unobserved locations.
The MSLE gap, by contrast, is small (0.0356 against 0.0385, $-7.5\%$).
Since MSLE here scores only above-mean traffic on a logarithmic scale,
this suggests that the spatial model's advantage lies mostly in
low-to-moderate activity regimes and in the meteorological channels,
which vary smoothly in space. At high-footfall peaks, where log-scale
errors dominate, the advantage is modest. Because the kriging RMSE and
MAE pool the traffic channel with two meteorological covariates, part
of the headline kriging advantage may reflect interpolation of spatially
smooth weather fields and not pedestrian dynamics. A traffic-only
kriging error in raw counts is needed to settle this.

With these qualifications, the strongest spatial evidence in the chapter
comes not from the baseline comparison but from Ablation~2
(Section~\ref{sec:exp3-abl-kriging}). There, the graph architecture is
held fixed and only the inductive supervision is removed. The resulting
collapse in kriging performance shows that the graph is necessary but
not sufficient for zero-shot transfer.

Relative to v3, the inductive evaluation is itself the principal
advance. The v3 simulated zero-history nodes produced pairwise
correlations of $r = 1.00$, which means the model assigned an identical
series to all five locations regardless of their coordinates. The 24 v4
inductive stations are real sensors with real ground truth, and the
existence of a meaningful gap between the spatially informed model and a
spatially blind reference is the first evidence in this thesis that RIoT
discriminates between unseen locations at all.
% AUTHOR-CHECK: a pairwise-correlation diagnostic across the 24 inductive
% stations' kriged series would allow a direct v3/v4 degeneracy comparison.

% ---------------------------------------------------------------------
\subsection{Deconstructing the architecture}
\label{sec:exp3-ablations}

\subsubsection{Joint reconstruction as a regulariser (Ablation 1)}
\label{sec:exp3-abl-impute}

Setting $w_{\text{impute}} = 0$ increases forecast RMSE by 7.6\% and MAE
by 9.7\%, and reduces ZIR by 17.6 percentage points. The
eightfold rise in imputation MSE (0.024 to 0.152) is expected: without
supervision, the reconstruction decoder is untrained, so this figure
tells us nothing about representational quality. The informative
quantity is the forecasting degradation. It shows that reconstruction
supervision improves the forecasting head even though the forecasting
head never sees its output.

The most parsimonious explanation is that the masked-reconstruction
objective acts as a regulariser on the shared temporal encoder. To
reconstruct artificially masked entries, the encoder must retain
information about within-window temporal structure and cross-channel
dependencies that a one-step-ahead forecasting loss alone would not
require. The resulting representation appears to be more robust to the
irregular observation patterns present at test time. This matches the
general finding from masked auto-encoding that reconstruction pretexts
encourage representations that transfer beyond the pretext task. It
extends that finding to a setting where the two tasks are trained
jointly and not sequentially. The kriging metrics also degrade modestly
(+5.9\% RMSE), consistent with a shared encoder whose general quality
has declined.

Two alternative explanations cannot be excluded from a single run. The
auxiliary loss may simply change the effective learning-rate schedule
of the shared parameters by adding gradient magnitude. It may also help
chiefly by teaching the encoder to handle the NaN-injected inputs it
encounters at test time, since artificial masking is applied during
evaluation as well as during training. The second explanation is not a
flaw, since robustness to missing input is the objective, but it implies
that the benefit is partly specific to the masking regime used here. On
balance, the evidence favours retaining joint reconstruction. It is not
an unnecessary auxiliary burden, and on this evidence it earns its
place in the objective.

\subsubsection{Inductive supervision and spatial generalisation (Ablation 2)}
\label{sec:exp3-abl-kriging}

Removing the kriging loss produces the most dramatic and most selective
effect in the ablation set. Forecast RMSE rises by only 3.2\% and
imputation MSE barely changes, but on the 24 inductive stations kriging
RMSE rises by 87\%, kriging MAE increases more than sixfold (0.052 to
0.332), and kriging MSLE more than doubles (0.0356 to 0.0794).

The most consequential observation is that, without the kriging loss,
the graph model performs \emph{worse} at inductive locations than the
non-spatial LSTM on all three kriging metrics (MAE 0.332 against 0.283;
MSLE 0.079 against 0.039). A spatial architecture that is trained only
on transductive forecasting and reconstruction does not produce usable
representations for unseen nodes. Its kriging decoder, deprived of
supervision, performs below a spatially blind reference. Message passing
over a learned graph provides the \emph{capacity} for spatial transfer,
but the inductive objective is what converts that capacity into
behaviour. Spatial generalisation in this architecture is therefore
learned, not emergent. For practical kriging, this means inductive
supervision must be part of the training objective. It should not be
treated as an optional refinement.

The small forecasting cost of removing the kriging term (+3.2\% RMSE)
indicates that the term also provides mild regularisation to the shared
representation. The simultaneous \emph{rise} in ZIR to 0.602 is
harder to interpret. One possibility is that the kriging term competes
with the zero-inflation head for shared capacity, so that freeing that
capacity lets the model fit the zero-filled targets more aggressively.
Section~\ref{sec:exp3-zero-recall} suggests why higher ZIR is not
necessarily desirable.

\subsubsection{The role of explicit missingness masks (Ablation 3)}
\label{sec:exp3-abl-masks}

Removing the two missingness-mask channels causes the largest
forecasting degradation of any ablation. RMSE rises by 9.6\%, MAE by
19.6\%, and ZIR falls by 25.9 points to 0.295. It is also the
only ablation for which we have validation figures showing a notably
larger generalisation gap. Validation-to-test RMSE rises by 7.6\% (18.50
to 19.91), compared with 1.2\% for the reference model and 3.0\% for the
batch-size-32 variant.

Two mechanisms, not mutually exclusive, explain why this ablation hurts
most. The first is \emph{disambiguation}. Once missing values are
zero-filled, the value channel alone cannot distinguish ``zero
pedestrians observed'' from ``no observation''. Both map to the same
standardised constant. The masks resolve that ambiguity for every step
of the input window. Without them, the encoder must infer outages
indirectly from the exact recurrence of a constant value, which is
unreliable in low-footfall periods where true zeros are common. The
larger validation-to-test gap fits this account: if the pattern of
outages shifts between chronological folds, a model that must learn
outage signatures implicitly will generalise less well than one that is
told about them explicitly.

The second mechanism is \emph{target alignment}, and it is the less
comfortable of the two. Because unobserved target intervals are scored
as $y = 0$, and because sensor outages persist over consecutive
timestamps, the mask state at the final input step is a strong predictor
of whether the target will be zero. A model with access to the masks can
learn to set $\hat{\pi} \to 1$ whenever a sensor is currently offline. It
is then rewarded twice: by the ZINB likelihood during training, and by
near-zero point error on those positions at test time. Part of the
accuracy lost in Ablation~3 is therefore accuracy on intervals that have
no physical ground truth. The ablation shows that the masks are
indispensable \emph{under the current target definition}. It does not
by itself show how much of their value comes from better modelling of
observed dynamics and how much comes from a better ability to predict
the zero-filling convention. We examine this directly in the next
section.

\subsubsection{The cost of permutation equivariance (Ablation 4)}
\label{sec:exp3-abl-pe}

Without sinusoidal positional encodings, the Transformer's temporal
self-attention is permutation-equivariant over the input window. The
cost is consistent but moderate: forecast RMSE and MAE rise by 4.5\% and
5.3\%, imputation MSE by 10\%, and the three kriging metrics by 4.5--6.5\%.

The mildness of this effect needs explaining, since urban pedestrian
flow is strongly ordered in time. Commuting peaks, lunchtime surges and
the evening decline all depend on sequence as well as on value. Two
factors probably limit the penalty. First, with a 12-step (3\,h) window,
much of the ordinal information is recoverable from content: a
monotonically rising count sequence during the morning implies its own
order, even to an order-agnostic mechanism. Second, other components
likely supply temporal position indirectly, such as any recurrent or
convolutional encoding stage and the forecasting head's read-out from
the final step.
% AUTHOR-CHECK: state whether time-of-day/day-of-week covariates or any
% order-sensitive layers precede the attention block.
Positional encoding is thus best described as a low-cost, reliable
refinement, not a load-bearing component. It removes a structural
inductive bias that runs against the physics of the problem, and it
buys roughly 0.8 RMSE for no measurable computational cost. We would
expect its contribution to grow with longer input windows, where
content-based recovery of order becomes harder.

% ---------------------------------------------------------------------
\subsection{ZIR and the risk of shortcut learning}
\label{sec:exp3-zero-recall}

The reference model reaches a ZIR of 0.554. Whether this counts
as a success depends entirely on what the zeros are.

\paragraph{The physical bound.}
Unobserved intervals account for 25.99\% of all timestamps, and 13.40\%
of all timestamps are unobserved \emph{and} fall between 00:00 and 05:00.
Night-time therefore accounts for $13.40 / 25.99 \approx 51.5\%$ of
unobserved intervals. Since these five hours make up only 20.8\% of the
day, roughly 64\% of night-time timestamps are unobserved, compared with
roughly 16\% of daytime timestamps. This fourfold concentration of
missingness at night is itself informative. It is consistent with a
missing-not-at-random mechanism in which sensors fail to report, or
report nothing, when there is no pedestrian activity. It supports
treating night-time gaps as plausible structural zeros. Daytime gaps are
much harder to justify as physical zeros. They are more plausibly sensor
or network outages during which pedestrians continued to pass.

It follows that, of the unobserved intervals that enter the ZIR
denominator as $y = 0$, only about half have a credible claim to being
true zeros. The 51.5\% figure is a heuristic, not a hard ceiling:
daytime zeros can be genuine (for example, during street closures), and
night-time gaps need not be. It does, however, provide a useful
reference point.

\paragraph{A plausibility band, not a target.}
A perfect zero identifier would not reach 100\% ZIR. It would flag
genuine observed zeros and night-time structural gaps, and it would
decline to flag daytime outages. We therefore treat a ZIR between
roughly 20\% and 80\% as physically plausible. Values below this band
would mean the zero-inflation component is effectively unused, while
values above it would mean the model is equating sensor failure with
absence of pedestrians. The 51.5\% night-time share falls near the
centre of the band. Every configuration in Table~\ref{tab:exp3-main}
(ZIR 0.295--0.602) lies within it. This is reassuring, since no
configuration has collapsed to either degenerate extreme. It also
means the band cannot, by itself, rank the configurations. Doing so
requires the stratified diagnostics set out at the end of this section.

\paragraph{Reading the reference model's ZIR.}
Because ZIR is computed over observed zeros and all zero-filled
intervals together, a ZIR of 0.554 cannot be decomposed without
stratification. Two scenarios bracket the interpretation.

\begin{table}[htbp]
    \centering
    \caption[Interpretive bounds on ZIR]%
    {Composition of the ZIR denominator and the interpretive
    reference provided by the night-time structural-zero heuristic.
    Shares are of all timestamps in the full record; the test-partition
    composition may differ.}
    \label{tab:exp3-zero-bound}
    \small
    \begin{tabular}{@{}lr@{}}
        \toprule
        \textbf{Quantity} & \textbf{Value} \\
        \midrule
        Unobserved timestamps (zero-filled, scored as $y=0$) & 25.99\% \\
        \quad of which night-time (00:00--05:00)              & 13.40\% \\
        \quad of which daytime                                & 12.59\% \\
        Night-time share of unobserved intervals              & 51.5\% \\
        Unobserved rate within night-time hours               & $\approx$64\% \\
        Unobserved rate within daytime hours                  & $\approx$16\% \\
        \midrule
        ZIR, full model                               & 55.4\% \\
        ZIR, no missingness masks                     & 29.5\% \\
        ZIR, no kriging loss (highest)                & 60.2\% \\
        \bottomrule
    \end{tabular}
\end{table}
% AUTHOR-CHECK: recompute the unobserved/night shares on the test
% partition specifically and add the observed-zero share.

If genuinely observed zeros are rare relative to zero-filled intervals,
the denominator is dominated by missing positions. In that case, a model
that fires $\hat{\pi} > 0.5$ on every night-time gap and on no daytime
gap would achieve a ZIR close to 0.515. The observed 0.554 exceeds
this, so the model must also be classifying some daytime gaps, or some
observed zeros, as zero. The configurations with even higher ZIR
(0.595 without positional encoding, 0.602 without the kriging loss) are
less accurate on every other metric. That pattern is more consistent
with a model predicting zeros too readily than with better structural
understanding. If observed zeros make up a larger share of the
denominator, the ZIR is more ambiguous, but the same logic applies to
the portion attributable to daytime gaps.

\paragraph{Evidence bearing on shortcut learning.}
Three observations together indicate that shortcut learning is not
merely a residual risk here. It is a feature that the current objective
actively encourages.

\begin{enumerate}
    \item \emph{The objective rewards it.} The ZINB likelihood is
    evaluated on zero-filled targets without a validity mask, and the
    mask channels expose the current outage state. The loss-minimising
    policy at an ongoing outage is therefore $\hat{\pi} \to 1$,
    irrespective of time of day. Nothing in the training signal
    distinguishes a night-time structural zero from a daytime outage.
    \item \emph{Ablation 3 is consistent with it.} Removing the masks
    roughly halves ZIR (0.554 to 0.295). A model that had learned
    structural zeros from temporal context (time of day, preceding
    decline in counts, location) would retain much of its ZIR without
    the masks. The size of the drop shows that the masks, not the
    dynamics, carry most of the zero-classification signal. Note that
    the two structural channels are a total-missingness mask and a
    \emph{night-restricted} missingness mask (missing \emph{and}
    00:00--05:00), and both are concatenated to the encoder input. The
    model is therefore given the night-time structural-zero heuristic
    directly, as a feature. Ablation~3 removes both channels together,
    and no run has yet disabled the night channel on its own, so the
    contribution of the explicit heuristic cannot be separated from
    that of general outage awareness.
    \item \emph{ZIR is dissociated from accuracy.} Across the
    eight configurations, higher ZIR is not associated with
    lower forecast error (Table~\ref{tab:exp3-delta}). The most accurate
    model by RMSE has one of the lowest ZIRs. If ZIR were tracking
    genuine structural knowledge, some positive association would be
    expected.
\end{enumerate}

None of this implies that the model has learned \emph{nothing} about
structural zeros. The strong MNAR concentration at night means that a
mask-driven zero prediction is correct more often than not during those
hours, and a model exploiting the mask will, as a by-product, be right
about many genuine structural zeros. The problem is one of
identifiability: under the present objective and metric, genuine
learning and shortcut exploitation produce the same numbers. We
therefore make no claim that a ZIR of 0.554 reflects learned
knowledge of structural zeros.

The same issue carries over to the forecast errors. Wherever the
reference model correctly sets $\hat{\pi} \to 1$ at a daytime outage,
it earns near-zero point error on a target that carries no physical
information. The headline RMSE and MAE are therefore partly measures of
agreement with the zero-filling convention. This is particularly
relevant when comparing configurations with different ZIR, and
when comparing against the LSTM baseline, which lacks the mechanism
entirely.

\paragraph{Diagnostics required.}
Resolving the ambiguity requires no new models, only re-scoring the
existing checkpoints under additional evaluation masks:
\begin{enumerate}
    \item report forecast RMSE and MAE restricted to \emph{observed}
    target positions, excluding zero-filled intervals;
    \item stratify ZIR into (a) observed zeros, (b) night-time
    unobserved and (c) daytime unobserved, and report precision and the
    calibration of $\hat{\pi}$ within each stratum;
    \item run a mask-perturbation test in which the final-step mask is
    flipped at inference, and measure the resulting change in
    $\hat{\pi}$. Large sensitivity to the mask with little sensitivity
    to the value history would be direct evidence of a shortcut.
\end{enumerate}
In the longer term, the principled remedy is to exclude unobserved
targets from the forecasting likelihood, or to model them explicitly
through a separate observation process. The zero-inflation component
should represent the physical process only, not the sensor process.

% ---------------------------------------------------------------------
\subsection{Optimisation dynamics: the small-batch result}
\label{sec:exp3-batch}

Reducing the batch size from 128 to 32 produces the lowest test RMSE
of any configuration (16.84, $-7.3\%$) and the lowest kriging RMSE
(0.213, $-9.3\%$). The improvement is visible on validation as well
(17.94 to 16.34), so it is not an artefact of selecting on the test
set. The gain is not uniform, however. Test MAE \emph{rises} by 4.7\%
(5.17 to 5.42), kriging MSLE rises by 5.1\%, and ZIR falls by
16.8 points.

This pattern of lower RMSE, higher MAE and lower ZIR points
to a redistribution of error, not a general improvement. RMSE is
dominated by the largest residuals, which in pedestrian data sit at
high-footfall peaks. MAE weights all residuals equally and is dominated
numerically by the large mass of low-count and zero positions. The
small-batch model has fitted the peaks better at a modest cost on the
low-count regime, including a reduced willingness to predict zeros. Its
RMSE/MAE ratio falls from 3.51 to 3.11, which confirms that the gain
comes from trimming the tail of large errors.

Several mechanisms could produce this.

\begin{itemize}
    \item \emph{Gradient noise as implicit regularisation.} Smaller
    batches raise the variance of stochastic gradient estimates, which
    is widely associated with convergence to flatter minima that
    generalise better. Peak-hour behaviour is rare and heterogeneous
    across days, and it is the regime most prone to overfitting, so it
    would plausibly benefit most.
    \item \emph{Greater gradient weight on rare events.} In a batch of
    128, a handful of peak-hour windows contribute a small, averaged
    share of each update. At a batch size of 32, individual batches that
    happen to contain peaks produce updates dominated by those peaks.
    This amounts to an implicit up-weighting of high-count samples in
    the ZINB mean component.
    \item \emph{More optimiser steps.} The two configurations are
    identical apart from batch size, with the same epoch cap (100),
    early-stopping patience (10 epochs) and learning rate
    ($3\times10^{-4}$, not rescaled). A batch size of 32 therefore
    performs four times as many parameter updates per epoch. Part of
    the gain may simply reflect a less under-trained model, not a
    qualitatively different solution.
\end{itemize}

The high-regularisation variant helps to separate these accounts.
Raising dropout from 0.1 to 0.3 \emph{worsens} RMSE by 7.5\% and MAE by
7.7\%, which suggests that the reference model is not overfitting in the
conventional sense and that additional explicit regularisation is
harmful. If the small-batch benefit were purely a matter of
regularisation, a similar direction of effect from dropout might have
been expected. That the two interventions move in opposite directions
favours the optimisation-trajectory accounts (more updates, greater
weight on rare events) over a pure flat-minima account. Both variants
also show a marked fall in ZIR (to 0.386 and 0.362). Any change
that perturbs the shared representation appears to weaken the
mask-driven zero policy first, which is a further indication that the
policy is fragile rather than deeply learned.

For downstream use, choosing between the two batch sizes is a choice of
loss function. If operational value lies in anticipating crowding peaks,
for example for public-safety or event management, the batch-size-32
model is preferable. If it lies in accurate low-activity and zero
forecasting, for example for energy-saving lighting schedules, the
reference model is preferable. We retain \texttt{wph1nhvj} as the
reference configuration for the ablation analysis, because all
ablations were run relative to it. Describing it as a performance
``ceiling'' would be inaccurate, however, since the batch-size result
shows that the v4 hyperparameter space has not been exhausted.

% ---------------------------------------------------------------------
\subsection{Limitations and threats to validity}
\label{sec:exp3-limitations}

\begin{enumerate}
    \item \textbf{Single-seed estimates.} Every configuration is a
    single run. The smaller effects reported here (kriging loss on
    forecasting, +3.2\%; positional encoding, +4.5\%) are of a size that
    seed variance in graph--Transformer models could plausibly produce.
    The larger effects (kriging loss on inductive kriging, masks on MAE
    and ZIR, the LSTM gap) are very unlikely to reverse, but
    their magnitudes should carry confidence intervals from at least
    three seeds.
    \item \textbf{Zero-filled targets.} As argued in
    Section~\ref{sec:exp3-zero-recall}, forecast errors and ZIR
    both include positions without physical ground truth. Every
    comparison in this chapter should be re-examined on observed-only
    targets.
    \item \textbf{Mixed node scope in forecast metrics.} Forecast RMSE,
    MAE and ZIR average the 57 transductive and 24 inductive stations
    together. They should be reported separately for each group.
    \item \textbf{Night-mask channel at inductive stations.} The
    night-restricted missingness channel is not blanked for held-out
    stations (Section~\ref{sec:exp3-protocol}). Kriging and
    inductive-station forecast metrics should be re-scored with it
    blanked.
    \item \textbf{Night channel not separately ablated.} Ablation~3
    removes the total and night-restricted mask channels together. An
    ablation that disables only the night channel is needed to establish
    how much of the zero-inflation behaviour comes from the explicitly
    supplied heuristic.
    \item \textbf{Mixed-channel kriging metrics.} Kriging RMSE and MAE
    pool traffic with meteorological channels in standardised space, so
    they do not isolate pedestrian kriging skill.
    \item \textbf{Confounded baseline.} The LSTM differs from RIoT in its
    distributional head, its access to mask channels and static
    covariates, and its spatial component. A
    graph-ablated RIoT (identity adjacency, all else fixed) would
    provide a clean estimate of the value of message passing.
    \item \textbf{Cross-generation comparability.} v3 and v4 differ in
    node population, temporal coverage and target handling. The
    comparison in Section~\ref{sec:exp3-v3-v4} supports a claim about
    evaluation reliability, not about architectural gain.
    \item \textbf{Single-factor design.} Ablations are one-at-a-time
    relative to a single reference. Interactions, such as whether the
    masks matter less once imputation supervision is present, are not
    identified.
\end{enumerate}

% ---------------------------------------------------------------------
\subsection{Summary}
\label{sec:exp3-summary}

Experiment~3 set out to test whether reconstruction-augmented,
graph-based forecasting could operate on heavily incomplete pedestrian
data without prior imputation, and to do so under an evaluation
protocol free of the leakage that affected earlier chapters. Five
conclusions follow.

\begin{enumerate}
    \item Under strict chronological splitting with a purge gap and
    training-only scaling, RIoT v4 achieves a test RMSE of 18.16 and a
    test MAE of 5.17, with a validation-to-test RMSE gap of only 1.2\%.
    The main contribution of the v4 overhaul is that these numbers can
    be trusted as out-of-sample estimates. They are not directly
    comparable with v3.
    \item Zero-shot spatial kriging to 24 unseen real stations is
    achieved and is clearly better than a spatially blind reference.
    It depends critically on explicit inductive supervision: without the
    kriging loss, the graph model kriges worse than a non-spatial LSTM.
    \item Joint masked reconstruction improves forecasting (7.6\% RMSE)
    and is best understood as a regulariser of the shared temporal
    encoder.
    \item Explicit missingness masks are the most influential single
    component. Part of their value, however, comes from enabling the
    model to reproduce the zero-filling of unobserved targets. Under the
    current objective and metric, genuine learning of structural zeros
    cannot be separated from shortcut exploitation of the mask, and the
    ZIR of 0.554 should not be read as evidence of the former.
    \item Hyperparameter changes shift the balance of error across
    regimes. Small-batch training lowers RMSE by 7.3\% by fitting
    footfall peaks better, but it raises MAE and weakens
    zero-classification.
\end{enumerate}

The most important next step is methodological, not architectural:
re-scoring the existing models on observed-only targets, with stratified
zero diagnostics. Until that is done, the forecasting results of this
chapter should be read as upper-bound estimates of performance on
physically observed pedestrian activity.
